%document
\documentclass[10pt]{beamer}
%theme
\usetheme{metropolis}
% packages
\usepackage{color}
\usepackage{listings}
\usepackage[main=ngerman, USenglish]{babel}
\usepackage[utf8]{inputenc}
\usepackage{multicol}
\usepackage{csquotes}


% color definitions
\definecolor{mygreen}{rgb}{0,0.6,0}
\definecolor{mygray}{rgb}{0.5,0.5,0.5}
\definecolor{mymauve}{rgb}{0.58,0,0.82}
\definecolor{paleorange}{HTML}{CC7832}

\lstset{
    backgroundcolor=\color{white},
    % choose the background color;
    % you must add \usepackage{color} or \usepackage{xcolor}
    basicstyle=\footnotesize\ttfamily,
    % the size of the fonts that are used for the code
    breakatwhitespace=false,
    % sets if automatic breaks should only happen at whitespace
    breaklines=true,                 % sets automatic line breaking
    captionpos=b,                    % sets the caption-position to bottom
    commentstyle=\color{mygreen},    % comment style
    % deletekeywords={...},
    % if you want to delete keywords from the given language
    extendedchars=true,
    % lets you use non-ASCII characters;
    % for 8-bits encodings only, does not work with UTF-8
    frame=single,                    % adds a frame around the code
    keepspaces=true,
    % keeps spaces in text,
    % useful for keeping indentation of code
    % (possibly needs columns=flexible)
    keywordstyle=\color{blue},       % keyword style
    % morekeywords={*,...},
    % if you want to add more keywords to the set
    numbers=left,
    % where to put the line-numbers; possible values are (none, left, right)
    numbersep=5pt,
    % how far the line-numbers are from the code
    numberstyle=\tiny\color{mygray},
    % the style that is used for the line-numbers
    rulecolor=\color{black},
    % if not set, the frame-color may be changed on line-breaks
    % within not-black text (e.g. comments (green here))
    stepnumber=1,
    % the step between two line-numbers.
    % If it's 1, each line will be numbered
    stringstyle=\color{mymauve},     % string literal style
    tabsize=4,                       % sets default tabsize to 4 spaces
    % show the filename of files included with \lstinputlisting;
    % also try caption instead of title
    language = Java,
	showspaces = false,
	showtabs = false,
	showstringspaces = false,
	escapechar = ,
        morecomment=[s][\textcolor{paleorange}]{@}{\ },
}

\def\ContinueLineNumber{\lstset{firstnumber=last}}
\def\StartLineAt#1{\lstset{firstnumber=#1}}
\let\numberLineAt\StartLineAt



\newcommand{\codeline}[1]{
        \alert{\texttt{#1}}
}

\setmonofont{Fira Code}

% This Document contains the information about this course.

% Authors of the slides
\author{Marcus Köhler}

% Name of the Course
\institute{Java-Kurs}

% Fancy Logo
\titlegraphic{\hfill\includegraphics[height=1.25cm]{../templates/fsr_logo_cropped}}

\usepackage{csquotes}
\usepackage{mathtools}

\title{Java}
\subtitle{Wiederholung}
\date{\today}

\begin{document}

\begin{frame}
    \titlepage
\end{frame}

\begin{frame}{Überblick}
    \setbeamertemplate{section in toc}[sections numbered]
    \tableofcontents
\end{frame}

\section{Wiederholung: Datentypen}
\begin{frame}{Quiz zu Datentypen}
    Welcher Datentyp ist ideal für die folgenden Werte? 
    \begin{center}
        \begin{tabular}{l l}
            \texttt{1337} & \pause \texttt{int} oder \texttt{short}\\
            \texttt{c} & \pause \texttt{char} \\
            \texttt{4.2} & \pause \texttt{float} oder \texttt{double} \\
            \texttt{12.345.678.910} & \pause \texttt{long} \\
            \texttt{Peter} & \pause \texttt{String}
        \end{tabular}
    \end{center}
\end{frame}

\section{Wiederholung: OOP}
\begin{frame}{Wiederholung: OOP}
    \center{\huge{Was macht ein Objekt aus?} \\
            \smallskip \huge{Wie passen Klassen dazu?}}
\end{frame}

\begin{frame}{Wiederholung:OOP}
    Ein \textit{Objekt} besteht aus den folgenden Teilen:
    \begin{itemize}[<+->]
        \item \textit{Attribute}, welche den momentanen Zustand und die Eigenschaften des Objekts darstellen. 
            \begin{itemize}
                \item Attribute sollten nicht für die Außenwelt direkt zugänglich sein \\ (-> Verkapselung). 
                \item Attribute können sowohl primitive Datentypen als auch andere Objekte sein.
            \end{itemize}
        \item \textit{Methoden}, die das sogenannte Verhalten des Objekts darstellen.
            \begin{itemize}
                \item Methoden regeln, was \& wie ein Objekt agiert.
                \item Die \textit{Signatur} einer Methode sieht folgendermaßen aus: \\
            \end{itemize}
            \onslide<+->{$
                        {\color{orange}\overbracket[1pt]{\texttt{\large{public static}}}^{\text{access\ modifiers}}\ }
                        {\color{blue}\underbracket[1pt]{\texttt{\large{void}}}_{\text{return\ type}}\ }
                        {\color{black}\overbracket[1pt]{\texttt{\large{main}}}^{\text{name}}\ }
                        {\color{gray}\underbracket[1pt]{\texttt{\large{(String[] args)}}}_{\text{parameter(s)}}}
                        $}
    \end{itemize}
\end{frame}

\begin{frame}{Wiederholung:OOP}
    Eine \textit{Klasse} dagegen ist sozusagen der Bauplan für ein Objekt.
    \begin{itemize}[<+->]
        \item Sie definiert, welche Attribute ein Objekt haben kann und welchen Typ sie haben.
        \item Ebenso definiert sie die Methoden, die ein Objekt verwenden kann, kann sie aber nicht \enquote{selbst} aufrufen.
        \item Außerdem definiert sie den \textit{Konstruktor}, eine spezielle Methode, die festlegt, 
            wie ein Objekt erzeugt wird und welche Anfangswerte es benötigt.
        \item Eine Klasse kann auch eigene Methoden und Variablen haben; diese sind mit dem Keyword \texttt{static} markiert.
    \end{itemize}
\end{frame}

\begin{frame}{OOP-Übung}
    -> Code
\end{frame}

\section{Wiederholung:Vererbung}
\begin{frame}{Wiederholung:Vererbung}
    \center{\huge{Was ist Vererbung? \\
            \smallskip Wozu braucht man Vererbung?}}
\end{frame}

\begin{frame}{Wiederholung:Vererbung}
    Vererbung wird verwendet, um gemeinsame Attribute und Methoden von mehreren Klassen in eine gemeinsame \textit{Superklasse} auszulagern. \\
    Dies hat die folgenden Auswirkungen:
    \begin{itemize}[<+->]
        \item Jede Instanz(d.h. ein Objekt) einer erbenden Klasse ist automatisch auch eine Instanz ihrer Superklasse.
        \item Attribute und Methoden, die in der Superklasse als \texttt{protected} oder \texttt{public} markiert sind, 
            sind auch in jeder erbenden Klasse (\enquote{Subklasse}) verfügbar.
        \item Wenn in der Superklasse ein Konstruktor definiert ist, \textit{muss} dieser im Konstruktor 
            der Subklasse mittels \texttt{super(\dots);} aufgerufen werden.
    \end{itemize}
\end{frame}

\begin{frame}{Übung:Vererbung}
    -> Code
\end{frame}

\section{Wiederholung:Interfaces}
\begin{frame}{Wiederholung:Interfaces}
    \center{\huge{Wofür braucht man Interfaces? \smallskip \\ Was macht ein Interface aus?}}
\end{frame}

\begin{frame}{Wiederholung:Interfaces}
    Ein \textit{Interface} ist sozusagen eine Vereinbarung, dass alle Objekte, die ein Interface implementieren, 
    mindestens die im Interface vorgegebenen Methoden umsetzt. Ein Interface ist durch die folgenden Eigenschaften gekennzeichnet:
    \begin{itemize}[<+->]
        \item Es definiert Methoden\textbf{signaturen}, die von den implementierenden Klassen \enquote{mit Funktionalität versehen werden müssen}.
        \item Es definiert Konstanten(d.h. Variablen mit dem Keyword \texttt{final}), die von allen Objekten der implementierenden Klassen geteilt werden.
        \item (Seit Java 8) Es \emph{kann} Methoden eine \enquote{Fallback}-Definition geben, falls eine Klasse sie nicht selbst definiert(mittels \texttt{default}).
    \end{itemize}
\end{frame}

\begin{frame}[fragile]{Wiederholung:Polymorphie}
    Es ist auch möglich, einer Referenz einer \enquote{übergeordneten} Klasse bzw. Interface mit einer \enquote{untergeordneten} Klasse zu belegen: 
    \begin{lstlisting}[gobble=8]
        public class Student extends Person {...}
        public class Video implements Streamable {...}
        [...]
        Person student = new Student(...);
        Streamable stream = new Video(...);
    \end{lstlisting}
    \textit{Hinweis:} Mit Polymorphie beschränkt man die \enquote{Fähigkeiten} des Objekts auf diejenigen, 
    die in der \enquote{übergeordneten} Klasse bzw. Interface definiert sind 
    und kann nicht mehr auf eventuelle zusätzlich definierte Methoden in den Unterklassen zugreifen.
\end{frame}

\begin{frame}{Übung:Interfaces \& Polymorphie}
    -> Code
\end{frame}

\end{document}
